\subsubsection{(1)}

Stefan--Boltzmanns lov gir følgende forhold:
\begin{align}
P = \varepsilon \sigma A T^4
\end{align}

der $P$ er effekten som str\aa{}les. Det er b\aa{}de str\aa{}ling \emph{ut} fra steinen:
\begin{align}
P_\text{ut} = \varepsilon \sigma A T^4
\end{align}

(der $T$ er temperaturen til steinen), og str\aa{}ling \emph{fra} luften \emph{til} steinen:
\begin{align}
P_\text{inn} = \varepsilon \sigma A T_\text{luft}^4
\end{align}

Vi kan legge disse sammen for \aa{} f\aa{} "netto" str\aa{}ling, alts\aa{} effekten som p\aa{}virker steinens indre energi:
\begin{align}
P_\text{netto} = \varepsilon \sigma A \left(T^4 - T_\text{luft}^4\right)
\end{align}

For \aa{} sette dette i sammenheng med temperaturen bruker vi følgende sammenheng:
\begin{align}
dU = c\, m\, dT
\end{align}

som sier at endring i indre energi, $dU$, er lik endring i temperatur, $dT$, ganget med $c\cdot m$, der $m$ er massen til steinen og $c$ er den spesifikke varmekapasiteten til granitt.

Vi deler p\aa{} $dt$ (en infinitesimal endring i tid) p\aa{} begge sider:
\begin{align}
\frac{dU}{dt} = c\, m\, \frac{dT}{dt}
\end{align}

Med andre ord: endring i indre energi over tid er lik endring i temperatur over tid ganget med $c\cdot m$.

Disse uttrykkene kan sl\aa{}s sammen. Endring i indre energi over tid er lik $-P_\text{netto}$ (fordi steinen er varmere enn omgivelsene og dermed taper energi):
\begin{align}
c\, m\, \frac{dT}{dt} = -P_\text{netto} = -\varepsilon \sigma A \left(T^4 - T_\text{luft}^4\right)
\end{align}

Vi løser for $\frac{dT}{dt}$ og f\aa{}r differensialligningen:
\begin{align}
\boxed{\;\frac{dT}{dt} = -\frac{\varepsilon \sigma A}{c\, m}\left(T^4 - T_\text{luft}^4\right)\;}
\end{align}

Siden luften har en konstant temperatur p\aa{} $T_\text{luft} = 20\,^\circ\text{C} = 293.15\,\text{K}$, kan vi sette inn tallverdien for $T_\text{luft}^4$:
\begin{align}
\frac{dT}{dt} = -\frac{\varepsilon \sigma A}{c\, m}\left(T^4 - 7.385 \cdot 10^{9}\,\text{K}^4\right)
\end{align}